% =============================================================================
% CHAPITRE 1 — INTRODUCTION GÉNÉRALE
% =============================================================================

\chapter{Introduction générale}

% =============================================================================
% 1.1 CONTEXTE GÉNÉRAL
% =============================================================================

\section{Contexte général}

L'eau constitue une ressource naturelle essentielle au développement des sociétés
et au fonctionnement de nombreuses activités humaines. Sa gestion efficace
représente aujourd'hui un enjeu majeur, particulièrement dans les régions
confrontées à des périodes de stress hydrique et à une augmentation constante
des besoins en consommation. Dans ce contexte, la réduction des pertes d'eau
et l'amélioration de la surveillance des installations hydrauliques constituent
des problématiques importantes pour les entreprises, les collectivités et les
gestionnaires d'infrastructures.

Les réseaux de distribution et les installations hydrauliques peuvent être
affectés par différents types de dysfonctionnements, notamment les fuites,
les consommations inhabituelles, les variations anormales de débit ou encore
les problèmes liés au fonctionnement des vannes. Lorsque ces événements ne
sont pas détectés rapidement, ils peuvent entraîner une consommation excessive
d'eau, une augmentation des coûts d'exploitation, des dégradations matérielles
et, dans certains cas, des conséquences environnementales importantes.

Les méthodes traditionnelles de surveillance reposent généralement sur des
contrôles périodiques et des inspections manuelles. Bien que ces approches
permettent d'identifier certains problèmes, elles présentent des limites
lorsqu'il s'agit de détecter rapidement des anomalies apparaissant entre deux
inspections. L'exploitation continue des données issues de capteurs connectés
constitue ainsi une approche permettant d'améliorer la surveillance des
installations et de réduire le délai de détection des événements anormaux.

L'intégration de l'Internet des Objets
(\textit{Internet of Things -- IoT}) permet justement de collecter en continu
des données provenant des équipements physiques. Ces données peuvent ensuite
être exploitées par des systèmes logiciels capables de les analyser
automatiquement et d'assister les opérateurs dans la supervision des
installations.

% =============================================================================
% 1.2 PRÉSENTATION DE L'ENTREPRISE ET DU PROJET
% =============================================================================

\section{Présentation de l'entreprise et du projet}

C'est dans ce contexte que s'inscrit le présent projet, réalisé au sein de
\textbf{SFM Connect}, branche spécialisée dans les solutions liées à
l'Internet des Objets (\textit{Internet of Things -- IoT}) du
\textbf{Groupe SFM}, à Tunis.

SFM Connect développe des solutions permettant de connecter des équipements
physiques à des plateformes numériques afin de collecter, traiter et exploiter
leurs données. Dans le domaine de la gestion de l'eau, le projet
\textbf{Droppy} s'inscrit dans cette démarche en proposant une solution
destinée à améliorer la supervision des installations hydrauliques à partir
des données provenant de capteurs IoT.

Ces capteurs transmettent régulièrement différentes informations relatives au
fonctionnement du réseau, notamment le débit d'eau
(\texttt{flow\_rate}), la consommation cumulative
(\texttt{consumption}) ainsi que l'état d'une vanne
(\texttt{state}), généralement représenté par les états
\texttt{ON} et \texttt{OFF}. Chaque équipement est identifié de manière
unique, notamment à travers son adresse MAC, permettant ainsi d'associer les
différentes mesures au capteur correspondant.

L'exploitation de ces données permet de passer d'une simple collecte
d'informations à une véritable approche de supervision intelligente. Le
système développé dans le cadre de ce projet vise ainsi à analyser
automatiquement les données de télémétrie afin d'identifier des comportements
inhabituels pouvant correspondre à des fuites ou à d'autres anomalies de
fonctionnement.

% =============================================================================
% 1.3 PRÉSENTATION DE LA SOLUTION DROPPY
% =============================================================================

\section{Présentation de la solution Droppy}

Le projet \textbf{Droppy -- Water Leak Detection \& Anomaly Analysis Suite}
consiste en la conception et la réalisation d'une plateforme logicielle
intelligente dédiée à la détection des fuites d'eau et à l'analyse des
anomalies provenant des capteurs IoT.

La solution adopte une architecture distribuée organisée autour de plusieurs
composants spécialisés. Elle associe une interface web développée avec
\textbf{Angular 19}, un backend principal développé avec
\textbf{Spring Boot 3} et un service dédié aux traitements d'intelligence
artificielle développé avec \textbf{FastAPI} et Python.

Cette organisation permet de séparer les différentes responsabilités du
système. Le backend constitue le service central de la plateforme et assure
notamment l'intégration des données, l'orchestration des traitements, la
gestion de la persistance ainsi que la communication avec l'interface
utilisateur. Le service d'intelligence artificielle prend en charge les
traitements nécessaires à l'analyse des données et à la détection des
comportements inhabituels.

Le tableau de bord constitue l'interface principale de supervision. Il permet
aux utilisateurs de consulter les informations relatives aux capteurs,
d'analyser les données historiques, de visualiser les anomalies détectées et
de recevoir des notifications lorsqu'un événement important est identifié.

Les choix techniques ainsi que l'organisation détaillée de cette architecture
seront présentés dans les chapitres consacrés à l'analyse des besoins et à la
conception de la solution.

% =============================================================================
% 1.4 PROBLÉMATIQUE
% =============================================================================

\section{Problématique}

La principale difficulté d'un système de surveillance des installations
hydrauliques réside dans la capacité à distinguer un comportement normal
d'un comportement réellement anormal. Une augmentation du débit peut, par
exemple, correspondre à une utilisation normale de l'installation, mais elle
peut également être le signe d'une fuite. De la même manière, un état
particulier de la vanne associé à un débit inattendu peut révéler une
incohérence dans le fonctionnement du système.

L'analyse d'un seul indicateur ne permet donc pas toujours d'obtenir une
détection suffisamment fiable. Il devient nécessaire de prendre en
considération plusieurs signaux provenant des capteurs et d'étudier leur
évolution dans le temps.

Par ailleurs, les règles statiques peuvent être efficaces pour identifier
certains cas connus, mais elles présentent des limites face à des
comportements inhabituels qui n'ont pas été explicitement définis à l'avance.
L'utilisation de techniques d'intelligence artificielle et d'apprentissage
automatique permet alors de compléter les règles métier en recherchant
automatiquement des comportements atypiques dans les données.

La problématique principale de ce projet peut ainsi être formulée comme suit :

\begin{quote}
\textit{
Comment concevoir et mettre en œuvre une plateforme intelligente, fiable
et évolutive permettant d'exploiter les données issues de capteurs IoT afin
de détecter automatiquement les fuites d'eau et les anomalies de
fonctionnement, tout en fournissant aux opérateurs une interface de
supervision capable de présenter les résultats et de générer des alertes
en temps réel ?
}
\end{quote}

% =============================================================================
% 1.5 APPROCHE PROPOSÉE
% =============================================================================

\section{Approche proposée}

Afin de répondre à cette problématique, le projet adopte une approche hybride
combinant des méthodes déterministes et des techniques d'apprentissage
automatique.

Dans un premier temps, les données provenant des capteurs sont nettoyées et
prétraitées afin de traiter les valeurs incohérentes et de préparer les
données pour les différentes étapes d'analyse. Une attention particulière est
accordée à l'exploitation des mesures de débit et de consommation en tant que
séries temporelles.

Dans un deuxième temps, des règles métier sont appliquées afin d'identifier
certaines situations particulières. Parmi ces mécanismes figure notamment
l'analyse du \textit{Minimum Night Flow} (MNF), permettant d'étudier les
consommations observées pendant les périodes nocturnes, ainsi que la détection
d'incohérences entre l'état d'une vanne et les mesures de débit.

En complément de ces règles, le système intègre des méthodes d'apprentissage
automatique. L'algorithme \textbf{Isolation Forest}, utilisé dans une approche
de détection non supervisée, permet d'identifier des observations présentant
un comportement inhabituel par rapport aux données observées.

Un mécanisme de prévision de la consommation est également utilisé afin de
comparer les tendances attendues aux comportements réellement observés.

Cette combinaison permet ainsi de dépasser les limites d'une détection
reposant uniquement sur des seuils fixes. Les résultats des différents
mécanismes d'analyse peuvent ensuite être exploités par le système de
supervision afin d'informer les opérateurs et de faciliter l'identification
des situations nécessitant une intervention.

% =============================================================================
% 1.6 ARCHITECTURE GÉNÉRALE DE LA SOLUTION
% =============================================================================

\section{Architecture générale de la solution}

L'architecture générale de Droppy repose sur trois composants principaux qui
collaborent pour assurer la collecte, le traitement, l'analyse et la
visualisation des données.

Le premier composant correspond au \textbf{frontend Angular 19}. Il constitue
l'interface utilisateur de la plateforme et permet notamment d'afficher les
informations relatives aux capteurs, les anomalies détectées, les données
historiques ainsi que différentes visualisations graphiques. Il permet
également de recevoir des événements en temps réel.

Le deuxième composant est le \textbf{backend Spring Boot 3}, qui constitue
le service central de l'application. Il assure l'intégration des différents
traitements, la gestion des données, la persistance des informations et la
communication avec l'interface utilisateur.

Le troisième composant correspond au \textbf{service d'intelligence
artificielle FastAPI}. Il regroupe les traitements liés à la préparation des
données, à la détection des anomalies et à la prévision. Il communique avec
le backend à travers des API afin de fournir les résultats nécessaires à la
plateforme.

Cette organisation permet de séparer clairement les responsabilités entre
l'interface utilisateur, la logique applicative et les traitements
d'intelligence artificielle. Elle facilite ainsi la maintenance du système
et permet d'envisager son évolution vers de nouveaux types de capteurs, de
nouvelles méthodes d'analyse ou de nouvelles fonctionnalités.

L'architecture détaillée, les flux de communication entre les composants
ainsi que les choix technologiques seront présentés dans les chapitres
consacrés à la conception de la solution.

% =============================================================================
% 1.7 OBJECTIFS DU PROJET
% =============================================================================

\section{Objectifs du projet}

L'objectif général de ce projet est de concevoir et développer une plateforme
intelligente permettant d'améliorer la détection des fuites d'eau et la
supervision des installations hydrauliques à partir de données IoT.

Plus précisément, les objectifs poursuivis sont les suivants :

\begin{itemize}

    \item Mettre en place un pipeline permettant de charger, nettoyer et
    prétraiter les données de télémétrie issues des capteurs IoT ;

    \item Développer des mécanismes de détection basés sur des règles métier
    afin d'identifier les comportements anormaux liés notamment au débit,
    à la consommation et à l'état des vannes ;

    \item Intégrer des méthodes d'apprentissage automatique, notamment
    l'algorithme \textbf{Isolation Forest}, afin de détecter automatiquement
    des comportements atypiques ;

    \item Mettre en place un mécanisme de prévision de la consommation
    permettant d'analyser les tendances et d'identifier d'éventuels écarts
    par rapport au comportement attendu ;

    \item Développer un service d'intelligence artificielle avec
    \textbf{FastAPI} afin d'exposer les fonctionnalités d'analyse,
    de prédiction et de prévision à travers des API REST ;

    \item Développer un backend central avec \textbf{Spring Boot 3}
    permettant d'assurer l'intégration des différents composants de la
    plateforme ;

    \item Mettre en œuvre un système de notifications en temps réel afin
    d'informer rapidement les opérateurs lorsqu'une anomalie importante
    est détectée ;

    \item Concevoir un tableau de bord interactif avec
    \textbf{Angular 19} permettant aux utilisateurs de visualiser les
    données, les anomalies et les informations relatives aux capteurs ;

    \item Assurer la persistance et la gestion structurée des données afin
    de permettre la consultation et l'analyse des historiques ;

    \item Mettre en place des tests permettant de vérifier le bon
    fonctionnement des différents composants de la solution.

\end{itemize}

% =============================================================================
% 1.8 MÉTHODOLOGIE ET ORGANISATION DU PROJET
% =============================================================================

\section{Méthodologie et organisation du projet}

La réalisation du projet a été organisée autour de plusieurs étapes
successives permettant de progresser progressivement de l'analyse du besoin
jusqu'à la validation de la solution.

La première étape a consisté à comprendre le fonctionnement du système IoT
et la structure des données provenant des capteurs. Cette phase a permis
d'identifier les différentes variables disponibles et de déterminer les
informations nécessaires à la détection des anomalies.

La deuxième étape a concerné le traitement et la préparation des données.
Les données brutes ont été analysées, nettoyées et transformées afin de
constituer une base exploitable pour les différents mécanismes de détection.

La troisième étape a consisté à concevoir et développer les mécanismes
d'analyse en combinant les règles métier et les méthodes d'apprentissage
automatique. Cette phase a notamment permis d'intégrer la détection des
comportements atypiques ainsi que les mécanismes de prévision de la
consommation.

La quatrième étape a été consacrée au développement de l'architecture
applicative. Les différents services ont été mis en place afin d'assurer
la communication entre les composants de la plateforme. En parallèle,
l'interface de supervision a été développée afin de permettre la
visualisation et l'exploitation des résultats.

Enfin, une phase de tests et de validation a permis de vérifier le
comportement des différentes composantes du système et d'évaluer la
cohérence des résultats produits.

% =============================================================================
% 1.9 STRUCTURE DU RAPPORT
% =============================================================================

\section{Structure du rapport}

Afin de présenter de manière progressive les différentes étapes de réalisation
du projet, ce rapport est organisé en plusieurs chapitres complémentaires.

Le premier chapitre présente le contexte général du projet, l'entreprise
d'accueil, la solution Droppy, la problématique ainsi que les objectifs
poursuivis.

Le deuxième chapitre est consacré à l'analyse des besoins et au cahier des
charges. Il présente notamment les acteurs du système, les besoins
fonctionnels et non fonctionnels, les contraintes ainsi que les critères
permettant d'évaluer la réussite de la solution.

Les chapitres suivants présentent progressivement l'analyse de l'existant,
la conception de la solution, l'architecture logicielle, les modèles UML
ainsi que les choix technologiques adoptés.

Une partie du rapport est ensuite consacrée à la réalisation de la plateforme.
Elle présente les principales fonctionnalités développées au niveau du
backend, du service d'intelligence artificielle et de l'interface de
supervision.

Les mécanismes d'intelligence artificielle et de détection des anomalies
sont également détaillés, notamment le prétraitement des données, les règles
métier, l'algorithme \textbf{Isolation Forest} ainsi que les mécanismes de
prévision.

Enfin, le rapport présente la stratégie de tests et de validation mise en
place afin d'évaluer le fonctionnement de la solution, avant de terminer
par une conclusion générale présentant les principales contributions
réalisées, les limites identifiées et les perspectives d'évolution de
la plateforme.

% =============================================================================
% CONCLUSION DU CHAPITRE
% =============================================================================

\section*{Conclusion du chapitre}

\addcontentsline{toc}{section}{Conclusion du chapitre}

Ce premier chapitre a permis de présenter le contexte général dans lequel
s'inscrit le projet \textbf{Droppy}, ainsi que les enjeux liés à la
surveillance intelligente des installations hydrauliques. L'analyse du
problème a mis en évidence les limites des approches traditionnelles et la
nécessité de disposer d'une solution capable d'exploiter en continu les
données issues des capteurs IoT.

La solution proposée repose sur une approche combinant traitement des
données, règles métier et techniques d'apprentissage automatique afin
d'améliorer la détection des comportements anormaux. L'architecture retenue
permet également de séparer les différentes responsabilités du système et
de fournir une interface de supervision adaptée aux besoins des utilisateurs.

Après avoir défini le contexte, la problématique et les objectifs du projet,
le chapitre suivant sera consacré à l'analyse détaillée des besoins et à
l'élaboration du cahier des charges. Cette étape permettra de transformer
les objectifs identifiés en exigences fonctionnelles et non fonctionnelles
précises, constituant ainsi la base de la conception de la solution.